% contex minimal starter document.
%
% Compile with:
%   docker compose run --rm compile-latex
% or, for a document tree elsewhere:
%   TEX_SRC_DIR=/path/to/doc TEX_FILE=paper.tex docker compose run --rm compile-latex
%
% The output lands in build/, never next to this file.

\documentclass[11pt, a4paper]{article}

% Page layout is a per-document decision; contex-cjk.sty deliberately does not
% set it for you.
\usepackage[a4paper, top=2.5cm, bottom=2.5cm, left=2cm, right=2cm]{geometry}

\usepackage{contex-cjk}   % options: [sans] or [serif] (default)
\usepackage{contex-stem}

\renewcommand{\baselinestretch}{1.3}

\title{\textbf{文件標題}}
\author{作者姓名}
\date{\today}

\begin{document}
\maketitle

\section{中英混排}
中文與 Latin text 混排時不需要任何手動標記：babel 依字元所屬的書寫系統自動切換字體與
斷行規則。行內數學如 $E = mc^{2}$ 也照常使用。

\section{數學}
行列式與方程式編號都由 \texttt{amsmath} 處理：
\begin{align}
  \int_{r}^{\infty} \frac{GMm}{r'^{2}} \, \mathrm{d}r' &= \frac{GMm}{r}. \label{eq:example}
\end{align}
式 \eqref{eq:example} 需要至少兩次編譯才會解析，因此本環境一律以 \texttt{latexmk} 驅動。

\section{表格與單位}
\begin{table}[htbp]
  \centering
  \begin{tabular}{lS[table-format=3.2]}
    \toprule
    項目 & {數值 (\si{\meter\per\second})} \\
    \midrule
    甲   & 12.30 \\
    乙   & 345.60 \\
    \bottomrule
  \end{tabular}
  \caption{以 \texttt{booktabs} 與 \texttt{siunitx} 排版的表格。}
\end{table}

\end{document}
